\section{The Emergence of Retrieval-Augmented Generation for Law (2020--2026)}
\label{sec:rag}

\subsection{From parametric knowledge to external legal evidence}

The development of large language models created a new possibility for legal AI: instead of
requiring a developer to encode every rule or training a separate model for every task, a
general-purpose model could interpret a question and generate a complete response in natural
language. However, this flexibility introduced a fundamental difficulty. The knowledge encoded in a
model's parameters is not the same as access to an authoritative, current, and verifiable legal
source.

This distinction is particularly important in law. Statutes may be amended, judicial decisions may
interpret or distinguish earlier authorities, and the applicability of a legal rule may depend on
jurisdiction and procedural context. A language model can generate a plausible explanation without
establishing which version of a provision applies or whether a cited judgment supports its
conclusion. The problem is therefore not simply that a model may lack information; it is that its
generated response may not provide sufficient evidence to verify its claims.

Retrieval-augmented generation (RAG) addresses this limitation by connecting a language model to an
external collection of documents. Rather than relying exclusively on its learned parameters, the
system retrieves relevant material and supplies that material to the generator as contextual
evidence. The generator then uses the retrieved information to formulate an answer. A basic legal
RAG pipeline has three stages:
\begin{enumerate}[leftmargin=1.4em, itemsep=2pt, topsep=2pt]
\item \textbf{Indexing:} Authoritative legal documents are collected, processed, divided into
retrievable units, and stored in a searchable representation.
\item \textbf{Retrieval:} A user's question is used to identify and rank potentially relevant
statutes, provisions, judgments, or passages.
\item \textbf{Evidence-conditioned generation:} The retrieved material is provided to the language
model, which generates a response using the supplied context.
\end{enumerate}

The distinction between parametric knowledge and external evidence is central. Model parameters
contain learned patterns and information acquired during training; the retrieval component provides
explicit documents that can be inspected at answer time. This separation can make a legal system
easier to update and its answers easier to verify. However, it does not guarantee correctness. A
generator may misunderstand retrieved evidence, combine incompatible authorities, or introduce claims
that the retrieved documents do not support.

RAG should therefore be understood as a mechanism for improving access to evidence, not as an
automatic guarantee of legal reliability.

\subsection{The legal RAG pipeline and its principal design choices}

The basic RAG architecture is relatively simple, but its effectiveness depends on several
interdependent decisions. The system must determine which documents to include, how to represent
their structure, how to retrieve relevant evidence, how to select the strongest passages, and how to
constrain the generator's response.

Legal documents make these decisions especially consequential. A statute is not merely a sequence of
independent paragraphs, and a judgment is not simply a bag of semantically related sentences. The
meaning of a passage can depend on its parent section, the facts of the case, the authority of the
court, and the relationship between the cited material and the legal issue.

For example, a query about a constitutional protection may retrieve a paragraph that describes the
general rule while missing a proviso or an associated procedural limitation. The retrieved paragraph
may be relevant but insufficient. If the generator treats it as complete evidence, it can produce an
answer that sounds reasonable yet omits a legally decisive qualification.

The quality of a legal RAG system is therefore determined by the complete pipeline, not merely by the
language model or embedding model. Table~\ref{tab:ragcomp} lists its principal components.

\begin{table*}[tp]
\centering
\caption{Principal components of a legal RAG architecture.}
\label{tab:ragcomp}
\footnotesize
\setlength{\tabcolsep}{5pt}
\renewcommand{\arraystretch}{1.3}
\begin{tabularx}{\textwidth}{>{\raggedright\arraybackslash}p{3.2cm} >{\raggedright\arraybackslash}X >{\raggedright\arraybackslash}X}
\toprule
\textbf{Component} & \textbf{Function} & \textbf{Main legal challenge} \\
\midrule
Legal corpus & Stores legislation, regulations, judgments, and other authoritative materials. & Establishing authority, jurisdiction, currency, and source provenance. \\
Document preprocessing & Extracts text and metadata from PDFs, HTML, and other formats. & Handling OCR errors, complex layouts, numbering, tables, and document boundaries. \\
Chunking and indexing & Divides documents into retrievable units and stores their representations. & Preserving relationships between provisions, clauses, exceptions, and parent sections. \\
Query processing & Interprets the user's information need and prepares it for retrieval. & Mapping everyday language to specialized legal terminology and exact statutory references. \\
Retrieval & Identifies and ranks candidate documents or passages. & Balancing exact legal expressions with conceptual relevance. \\
Reranking and evidence selection & Reorders candidates and selects material for the generator. & Distinguishing merely relevant passages from legally sufficient evidence. \\
Generation & Produces an answer using the query and retrieved context. & Avoiding unsupported claims and interpreting evidence within its proper context. \\
Verification & Checks claims and citations against the retrieved sources. & Determining whether each citation actually supports the associated claim. \\
\bottomrule
\end{tabularx}
\end{table*}

This architecture also clarifies why retrieval quality and generation quality must be evaluated
separately. A retrieval system may identify the correct provision, but the generator may misinterpret
it. Conversely, a capable generator may be unable to produce a reliable answer because the retrieval
component omitted the decisive evidence.

\subsection{From basic RAG to more adaptive systems}

Legal RAG systems can be understood as a progression in how they manage retrieval and generation.
Basic systems execute a relatively fixed retrieval workflow. More advanced systems add query
transformation, reranking, filtering, or iterative retrieval. Modular systems separate these
functions into replaceable components, while agentic systems use language-model-driven planning or
tool selection to coordinate multiple steps.

These categories are useful analytical distinctions, although individual systems may combine several
of them rather than fitting neatly into a single category. Table~\ref{tab:ragevol} summarizes the
progression.

\begin{table*}[tp]
\centering
\caption{Evolution of legal RAG architectures.}
\label{tab:ragevol}
\footnotesize
\setlength{\tabcolsep}{5pt}
\renewcommand{\arraystretch}{1.3}
\begin{tabularx}{\textwidth}{>{\raggedright\arraybackslash}p{2.6cm} >{\raggedright\arraybackslash}X >{\raggedright\arraybackslash}X >{\raggedright\arraybackslash}X}
\toprule
\textbf{Architecture} & \textbf{Operating principle} & \textbf{Potential benefit} & \textbf{Main limitation} \\
\midrule
Basic or naive RAG & Retrieve passages for a query and provide them to a generator. & Simple, interpretable baseline with explicit external context. & Sensitive to chunking, retrieval failures, and irrelevant context. \\
Advanced RAG & Improve preprocessing, query handling, retrieval, reranking, or context selection. & Can improve evidence relevance and reduce unnecessary context. & More components introduce additional parameters, complexity, and evaluation requirements. \\
Modular RAG & Organize retrieval, transformation, ranking, and generation as separable components. & Supports component-level testing and replacement. & Improvements depend on the quality of component integration. \\
Agentic RAG & Use an agent or multiple agents to plan, select tools, refine queries, or retrieve evidence iteratively. & Can support complex questions requiring multiple retrieval steps. & May increase latency, cost, and error propagation; additional agent steps do not guarantee better answers. \\
\bottomrule
\end{tabularx}
\end{table*}

These categories should not be interpreted as a universal maturity ranking. A carefully designed
basic RAG system may outperform an unnecessarily complex agentic workflow on a narrow legal
question-answering task. The relevant question is whether a more elaborate architecture produces
measurable improvements in evidence quality, answer reliability, or operational efficiency.

The distinction also matters for evaluation. An iterative system may retrieve additional evidence
after an initial answer, but its performance should be measured against a comparable baseline.
Without controlled evaluation, it is difficult to determine whether improvements arise from query
refinement, better retrieval, a stronger generator, or simply additional opportunities to produce an
answer.

\subsection{Representative legal RAG systems}

Recent studies illustrate how retrieval-augmented generation is being adapted to different legal
tasks. They also show why legal RAG should not be treated as a single architecture with one universal
evaluation method.

\citet{buffa2025}, in \emph{Enhancing Legal Document Building with Retrieval-Augmented Generation},
present JusBuild, a system designed to assist legal professionals in drafting structured legal
documents. The architecture segments existing documents into functional sections, stores their
vector representations, retrieves relevant precedent sections, and provides additional AI-generated
suggestions. An important design choice is its human-in-the-loop approach: the system assists with
drafting while leaving the final formulation under the professional user's control. The study
evaluates the architecture using datasets that differ in language, document template, and judicial
matter. Its focus is document construction and drafting assistance rather than a general legal
question-answering benchmark.

\citet{ren2026} address a different problem in \emph{Automated Data Synthesis and Retrieval-Augmented
Generation for Legal Large Language Models}. Their framework combines synthetic legal data generation
with quality evaluation and optimization, including an LLM-based judge and an enhanced Elo-ranking
mechanism. They evaluate the approach using LawBench, a benchmark for Chinese legal language models
that covers memory, comprehension, and application tasks. Their contribution is therefore not limited
to query-time retrieval: it also addresses the preparation and quality of training data for legal
models.

These two studies illustrate an important distinction. JusBuild uses retrieval to provide relevant
precedent sections during document drafting. Ren et al.\ combine synthetic data preparation, model
adaptation, and RAG-oriented evaluation to improve legal LLM performance. Their objectives,
architectures, and evaluation settings differ, so their reported results should not be interpreted as
a direct head-to-head comparison.

A further development is the explicit evaluation of the retrieval component itself. Pipitone and Houir
Alami's LegalBench-RAG \citep{pipitone2024} was introduced to assess retrieval for legal RAG,
emphasizing the retrieval of precise text segments rather than merely identifying relevant documents. This addresses a critical
evaluation gap: a system may retrieve the correct document but fail to identify the specific passage
needed to support an answer.

Together, these studies demonstrate that legal RAG research addresses several related but distinct
questions: how to retrieve relevant evidence, how to support document drafting, how to improve legal
models through data synthesis, and how to evaluate retrieval independently from generation.

\subsection{Agentic legal AI: from retrieving evidence to coordinating legal tasks}

The next development is the integration of agentic workflows into legal RAG. A conventional RAG
pipeline usually follows a predefined sequence: retrieve evidence, construct a prompt, and generate
an answer. An agentic system can instead use a language model to plan subtasks, choose tools,
reformulate queries, inspect intermediate results, and decide whether additional retrieval is
necessary.

This flexibility is relevant to legal questions that cannot be answered adequately by a single
retrieval operation. Consider a question involving a statutory provision and a judicial
interpretation. A system might need to identify the relevant statute, retrieve judgments interpreting
it, determine whether those judgments concern comparable facts, and then synthesize the evidence. An
agentic workflow can coordinate these operations as separate steps.

However, additional reasoning steps also create additional opportunities for error. An agent may
begin with an incorrect assumption, select an unsuitable tool, retrieve a misleading authority, or
propagate an unsupported conclusion into subsequent stages. Multiple agents do not automatically
provide independent verification, especially when they share the same underlying model or evidence.

For legal applications, the important design question is therefore not how many agents a system
contains, but whether its workflow improves evidence quality and makes the resulting reasoning more
verifiable. Each consequential claim should remain traceable to relevant source material, and the
system should be able to identify when its evidence is incomplete or contradictory.
Table~\ref{tab:agentic} contrasts conventional and agentic legal RAG.

\begin{table*}[tp]
\centering
\caption{Conventional and agentic legal RAG.}
\label{tab:agentic}
\footnotesize
\setlength{\tabcolsep}{5pt}
\renewcommand{\arraystretch}{1.3}
\begin{tabularx}{\textwidth}{>{\raggedright\arraybackslash}p{2.8cm} >{\raggedright\arraybackslash}X >{\raggedright\arraybackslash}X}
\toprule
\textbf{Dimension} & \textbf{Conventional RAG} & \textbf{Agentic RAG} \\
\midrule
Workflow & Predetermined retrieval and generation sequence. & Workflow may adapt to intermediate results. \\
Query handling & Usually retrieves from the original or a transformed query. & Can plan and execute multiple related retrieval operations. \\
Evidence gathering & Generally performed within a fixed retrieval step. & Can iteratively gather and compare additional evidence. \\
Error handling & Depends on predefined fallback and verification mechanisms. & Can attempt recovery, but may also propagate errors across steps. \\
Evaluation priority & Retrieval effectiveness, grounded generation, latency, and cost. & The same metrics, plus planning quality, tool selection, and error propagation. \\
\bottomrule
\end{tabularx}
\end{table*}

The broader implication is that agentic orchestration should complement, rather than replace,
reliable retrieval and evidence verification. A complex workflow built on incomplete or poorly
structured legal evidence remains unreliable.

\subsection{Synthesis: what legal RAG solves, and what it leaves open}

Retrieval-augmented generation represents an important shift in legal AI because it provides a
mechanism for connecting language generation to explicit legal sources. It can make answers more
specific to a selected corpus and provide material against which claims can be checked. Its modular
structure also permits retrieval, generation, and verification to be improved independently.

Yet RAG does not eliminate the underlying difficulties of legal information processing. Its output
remains dependent on the quality of the corpus, the preservation of document structure, the relevance
of retrieved passages, and the generator's use of that evidence. Nor does a citation automatically
establish correctness: the cited passage must support the claim, and the claim must be interpreted in
the appropriate legal context.

This leads to the central architectural problem addressed in the next section: how can a retrieval
system identify the right legal evidence while preserving the structural and semantic context
required to interpret it correctly?
